\chapter{Yield Function Customization and Constitutive Curve Visualisation}

\section{Why Introduce an Explicit Yield Function?}

Up to this point the plasticity architecture already had three explicit
compile-time axes:

\begin{itemize}
  \item \texttt{YieldCriterion}: geometry of the admissible surface in stress
        space;
  \item \texttt{HardeningLaw}: evolution of the admissible radius or position;
  \item \texttt{FlowRule}: direction of irreversible strain increment.
\end{itemize}

That is already a good decomposition, but semantically one piece was still
implicit: the scalar yield function value
\[
f(\bsigma, \alpha)
\]
that decides whether the trial state is admissible or not.

In the original implementation, that scalar was assembled directly inside
\texttt{PlasticityRelation} as
\[
f = q(\bsigma) - \sigma_y(\alpha).
\]
This is mathematically standard, but architecturally it hides a real
customization point.

\section{Criterion vs Function}

The distinction matters:

\begin{itemize}
  \item a \emph{yield criterion} describes the shape of the surface, for
        example von Mises or Drucker--Prager;
  \item a \emph{yield function} is the scalar map that evaluates the current
        trial state against that surface and the evolving hardening state.
\end{itemize}

For classical associative plasticity they are closely related, but they are
not identical concepts.  The criterion gives the stress measure
\(q(\bsigma)\) and its gradient; the yield function combines that measure with
the hardening law into a scalar overstress.

\section{New Compile-Time Contract}

The new concept is
\texttt{YieldFunctionPolicy<F, N, Y, H>}.  It requires

\begin{lstlisting}[language=C++]
f.value(yield, trial, hardening, state) -> double
\end{lstlisting}

where:

\begin{itemize}
  \item \texttt{yield} is the \texttt{YieldCriterion};
  \item \texttt{trial} is the effective trial state;
  \item \texttt{hardening} is the \texttt{HardeningLaw};
  \item \texttt{state} is the hardening-law state.
\end{itemize}

The default implementation is \texttt{StandardYieldFunction}, which reproduces
the previous behaviour exactly:
\[
f = q(\texttt{trial}) - \sigma_y(\texttt{state}).
\]

\section{How \texttt{PlasticityRelation} Uses It}

\texttt{PlasticityRelation} is now parameterised as

\begin{lstlisting}[language=C++]
PlasticityRelation<Policy,
                   YieldCriterion,
                   HardeningLaw,
                   FlowRule,
                   YieldFunction>
\end{lstlisting}

with \texttt{YieldFunction = StandardYieldFunction} by default.

This means the architectural axes are now:

\begin{enumerate}
  \item measure space / constitutive policy;
  \item yield-surface geometry;
  \item hardening evolution;
  \item flow direction;
  \item scalar yield-function evaluation.
\end{enumerate}

\section{Important Current Limitation}

The current return-mapping algorithm is still the standard backward-Euler
radial-return form.  Therefore the new \texttt{YieldFunction} policy is
architecturally explicit, but algorithmically it is still expected to remain
compatible with the same correction structure.

In practice, this means the present implementation is best suited to yield
functions that preserve the same derivative structure as the classical
overstress map, for example shifted or scaled variants that do not change the
underlying flow/tangent derivation qualitatively.

This is not a weakness of the concept itself; it is simply an honest statement
about the current integrator.  If later a different local algorithm is added,
the \texttt{YieldFunction} axis is already explicit and can be reused.

\section{Why This Separation Helps}

The gain is mostly semantic today, but it matters for future work:

\begin{itemize}
  \item classical \(f = q - \sigma_y\);
  \item capped or shifted overstress maps;
  \item pressure-sensitive surfaces with the same hardening but different
        scalar evaluation;
  \item dynamic/viscoplastic extensions where the overstress enters a rate law.
\end{itemize}

Separating the scalar yield function early is cheap and removes one more
implicit coupling from the plasticity kernel.

\section{Constitutive Protocol Sampling}

The new header \texttt{ConstitutiveProtocol.hh} introduces a tiny solver-free
utility for material tests:

\begin{lstlisting}[language=C++]
sample_constitutive_protocol(site, history,
                             abscissa_of, ordinate_of)
\end{lstlisting}

It applies a prescribed history of kinematic states and records:

\begin{itemize}
  \item abscissa value, typically strain or generalized strain;
  \item ordinate value, typically stress or generalized force;
  \item path parameter (currently the step index);
  \item discrete step number.
\end{itemize}

The commit path is resolved at compile time:

\begin{itemize}
  \item if the object exposes \texttt{commit(k)}, it is used;
  \item otherwise, if it exposes \texttt{update(k)}, that path is used.
\end{itemize}

This lets the same utility work for:

\begin{itemize}
  \item raw constitutive relations with embedded state;
  \item typed constitutive sites \texttt{MaterialInstance<...>};
  \item type-erased \texttt{Material<Policy>} handles.
\end{itemize}

\section{Why the Protocol Is Strain-Driven}

The present utility is displacement/strain-controlled on purpose.  Local
constitutive relations are naturally evaluated as
\[
\bvarepsilon \mapsto \bsigma.
\]
Force-controlled material tests would require solving a local inverse problem
at each step, which belongs to a different abstraction level.  Nothing in the
current design prevents that later; it is simply not mixed into the present
utility.

\section{VTK Representation of Material Curves}

The new \texttt{VTKConstitutiveCurveWriter} exports each curve as a
\texttt{vtkPolyData} polyline in the geometric plane
\[
(x,y,z) = (\text{strain}, \text{stress}, 0).
\]

This is intentionally simple and robust:

\begin{itemize}
  \item ParaView can render the loop directly as geometry;
  \item the line can be coloured by stress, strain, or step;
  \item the same file can be inspected either as a curve in space or as a
        source for further filters.
\end{itemize}

For each point, the writer stores arrays:

\begin{itemize}
  \item \texttt{strain} (or the chosen abscissa name);
  \item \texttt{stress} (or the chosen ordinate name);
  \item \texttt{step};
  \item \texttt{path\_parameter}.
\end{itemize}

Multiple curves are exported as a \texttt{.vtm} multiblock collection, with one
block per material/protocol.

\section{Curves Generated by the New Test}

The regression suite now writes:

\begin{itemize}
  \item \texttt{material\_hysteresis\_curves\_menegotto\_cyclic.vtp}
  \item \texttt{material\_hysteresis\_curves\_kent\_park\_cyclic.vtp}
  \item \texttt{material\_hysteresis\_curves\_j2\_backbone.vtp}
  \item \texttt{material\_hysteresis\_curves.vtm}
\end{itemize}

The first two show true hysteretic behaviour:

\begin{itemize}
  \item steel reversal with Menegotto--Pinto smooth transition and
        Bauschinger-type response;
  \item concrete compression/unloading/reloading with cracking and tension
        cut-off in Kent--Park.
\end{itemize}

The J2 block is kept as a backbone-style reference for the current local
plasticity implementation.

\section{What the New Tests Actually Verify}

Two kinds of behaviour are now checked:

\begin{enumerate}
  \item \textbf{Architectural composition}: a custom compile-time
        \texttt{YieldFunction} changes stress and tangent while the rest of the
        plasticity composition remains unchanged.
  \item \textbf{Post-processing contract}: the VTK writer produces valid
        polyline datasets with named arrays and block sidecars that can be read
        back and inspected.
\end{enumerate}

This is important because it validates not only mathematics, but also the
intended user workflow for constitutive debugging and calibration.

\section{Relevance for Reinforced Concrete}

For reinforced concrete across scales, graphical constitutive inspection is not
optional.  The steel and concrete laws often look acceptable numerically while
still behaving incorrectly under reversal, unloading, cracking, confinement, or
stiffness recovery.  The new protocol sampler and VTK curve writer give a
direct way to inspect that behaviour from the same codebase that later runs the
structural analysis.

This is especially useful when calibrating:

\begin{itemize}
  \item reinforcement hysteresis;
  \item confined versus unconfined concrete;
  \item fiber-section assembly;
  \item future damage or fracture laws in continuum concrete.
\end{itemize}

\section{Strategic Next Step}

The next logical move is to reuse this protocol infrastructure for more general
constitutive diagnostics:

\begin{itemize}
  \item generalized force--deformation curves for beam and shell sections;
  \item cyclic protocols for section reductions and fiber sections;
  \item eventually, local traction--separation laws and continuum damage /
        fracture models.
\end{itemize}

That would extend the same workflow from uniaxial laws to reduced structural
models and later to continuum inelasticity.
